\subsection{Laby}
Počítačovým laboratořím říká každý správný matfyzák pouze \textit{lab} a na
Matfyzu jich je naštěstí relativně dost. V každé budově je nejméně jeden takový,
do něhož mají přístup normální uživatelé (studenti Matfyzu).
Kromě toho existuje ještě spousta katedrových počítačů, u kterých sedávají
studenti vyšších ročníků.

K přístupu do labu potřebujete průkaz studenta a uživatelské konto nebo přístup k SISu.
Potřebujete-li založit konto, procedura bývá taková, že si člověk zažádá u
služby v daném labu, žádosti je zpravidla vyhověno a často okamžitě je konto
založeno. V Troji a na Karlově má přístup k počítačům automaticky každý student
univerzity. Přihlásíte se tam pomocí stejných údajů jako do SISu. Někdy je pro
založení účtu potřeba odchytit správce a absolvovat školení o specifikách daného
labu.

Obsazenost labů závisí na denní době, fázi semestru a poloze v Praze. Nejplnější
bývají laby v zimním semestru v období psaní zápočtových programů. Nejvíce místa
je obvykle na Karlově, jelikož zdejší dav není lehké najít. Nezapomeňte se řídit řádem daného labu a
směrnicí děkana č. 4/2008, ze které např. vyplývá, že se máte k počítačům chovat
slušně a že nesmíte používat nelegální software ani psát vulgární e-maily.

V každém labu je tiskárna, tisk stojí jednotně jedna koruna za černobílou stranu
A4 a 5 až 15 korun za barevnou. Ne všude je barevná tiskárna. Ve většině labů
najdete i scanner.

Každý lab má svoji webovou stránku, rozcestník najdete na: \url{https://www.mff.cuni.cz/cs/vnitrni-zalezitosti/it-a-sluzby/pocitacove-laboratore}.

Pro více informací o labech se podívejte do elektronické Kuchařky - \url{http://www.matfyzak.cz/kucharka/index.php?title=Laby}.



\subsubsection{Alternativy}
Situace, kdy má počítačová laboratoř zavřeno a spolubydlící si zabral buňku pro
radostné trávení času s přítelkyní, jsou časté a může se stát, že budou
kolidovat s chvílemi, kdy opravdu potřebujete pracovat. Máte-li notebook, můžete
využít noční studovnu v Národní technické knihovně v Dejvicích. Je otevřená
kdykoliv, kdy není otevřená vlastní knihovna, je v ní k dispozici WiFi,
elektrický proud, záchod a automaty s pitím a bagetami. Registrovat se ale je
potřeba někdy v běžných otevíracích dobách – stojí to padesát korun na rok.